% Summary: Working on the FLAIR Desktop (AN-14 summary tables, "To do this /
% Take these steps"; the green Summary pages of the 1991 Macintosh User's
% Guide, apple-macintosh.md Sec 5.2, 9).  Every step restates a task taught
% in Part I, against the same sources and gates.
% GATES: test-flair-disk-windows test-flair-file-ops test-flair-finder-cmds test-flair-app-launch test-flair-itext test-flair-empty-trash test-flair-shutdown test-flair-exit-to-dos test-pubs-flair-screens
\Gates{test-flair-disk-windows, test-flair-file-ops, test-flair-finder-cmds,
  test-flair-itext, test-flair-empty-trash,
  test-flair-shutdown, test-flair-exit-to-dos, test-pubs-flair-screens}
\BackChapter{Summary: Working on the FLAIR Desktop}
\index{Summary}

\section{Windows}
\begin{SummaryTable}
\SummaryRow{Open the data disk}{Double-click the \texttt{INITECH} icon.}
\SummaryRow{Open a folder}{Double-click the folder, or select it and press
\key{Ctrl-O}.}
\SummaryRow{Make a window active}{Click anywhere in it.}
\SummaryRow{Move a window}{Drag its title bar.}
\SummaryRow{Change a window's size}{Drag its size box, at the lower right
corner.}
\SummaryRow{Fill the screen with a window}{Click its zoom box.  Click again
to restore it.}
\SummaryRow{Roll a window up}{Click its collapse box.  Click again to roll
it down.}
\SummaryRow{See more of a window's contents}{Click the scroll arrows, click
the grey area of a scroll bar, or drag the scroll box.}
\SummaryRow{Close a window}{Click its close box.  For the front Finder
window you can press \key{Ctrl-W}.}
\end{SummaryTable}

\section{Files and folders}
\begin{SummaryTable}
\SummaryRow{Select an icon}{Click it.}
\SummaryRow{Select several icons}{Click the first, then \key{Shift}-click
the others; or drag a rectangle round them; or press \key{Ctrl-A} to select
them all.}
\SummaryRow{Make a folder}{Click the window of the folder it should be in,
then press \key{Ctrl-N}.}
\SummaryRow{Move a file into a folder}{Drag its icon onto the folder icon,
or into the folder's window.}
\SummaryRow{Move a file to the top level of the disk}{Drag its icon onto the
\texttt{INITECH} icon.}
\SummaryRow{Tidy a window}{Choose \menuitem{Clean Up} or \menuitem{Arrange
(by Name)} from the \menuname{View} menu.}
\end{SummaryTable}

\section{Applications and documents}
\begin{SummaryTable}
\SummaryRow{Open an application}{Double-click its icon in the \file{APPS}
folder.  For the Calculator, choose it from the Apple menu.}
\SummaryRow{Change to another application}{Click one of its windows.  Click
the desktop for the Finder.}
\SummaryRow{Save a document}{Press \key{Ctrl-S}.  The first time, type a
name and press \key{Enter}.}
\SummaryRow{Open a document}{In its application, press \key{Ctrl-O},
select the document and press \key{Enter}.}
\SummaryRow{Quit an application}{Press \key{Ctrl-Q}.}
\end{SummaryTable}

\section{Throwing things away}
\begin{SummaryTable}
\SummaryRow{Throw away a file or folder}{Drag its icon onto the Trash.}
\SummaryRow{Take something out of the Trash}{Double-click the Trash, then
drag the item out of the Trash window.}
\SummaryRow{Delete what is in the Trash}{Choose \menuitem{Empty Trash} from
the \menuname{Special} menu, then click \texttt{OK}.}
\end{SummaryTable}

\section{Finishing work}
\begin{SummaryTable}
\SummaryRow{Turn the computer off}{Click the desktop, choose \menuitem{Shut
Down} from the \menuname{Special} menu, then turn the computer off.}
\SummaryRow{Restart the computer}{Choose \menuitem{Restart} from the
\menuname{Special} menu.}
\SummaryRow{Go to the InitechDOS prompt}{Choose \menuitem{Exit to
InitechDOS} from the \menuname{Special} menu, click \texttt{Exit}, and press
\key{Enter} twice.}
\end{SummaryTable}
